\chapter*{Abstract}
\thispagestyle{empty}

For the ASCEND research group at the École de technologie supérieure, evidence fragmentation during the early lifecycle phases of nanosatellite electrical power systems hinders the traceability between needs and requirements, architecture decisions, and the analytical evidence that supports them. Furthermore, it hinfers the reproducibility of analyses and the capability to detect design inconsistencies before advacing to physical verification and validation phases.
This project mitigates this problem by developing a model-based systems engineering framework oriented at the early justification of nanosatellite electrical power systems through traceable and reproducible evidence.
The framework is implemented using Python and Github, at a positive economic value, and fulfills requirements such as retaining parametric requirements, identifying impacts of a proposed mission, requirement, component, or model change, proposing COTS EPS components from available requirements, and recommending and executing a set of EPS analyses and verification methods.
Two case studies compared the analysis outputs of the framework and those of BIRDS-4 and OreSat publications under the same mission and EPS inputs, and found consistent values and findings.
User trials showed a significant decrease in the time taken to perform controlled early justification tasks for stakeholders leveraging the framework.
The framework should be extended to other nanosatellite subsystems, and the release of the entire codebase as a public repository enables this.
\bigskip

\textbf{Keywords:} \scriptKeywords

\cleardoublepage

%%% Local Variables:
%%% mode: latex
%%% TeX-master: "main"
%%% End:
